% !TeX root = ../../Book1.tex
% 纲要标题：### 6.4 标准模型的约化
% 纲要位置：几何学计划纲要.md，第 342 行。
% 正文状态：初稿；尚未正式审核。

% 纲要内容开始
% 
% * 手柄和交叉帽的抽取
% * 可定向标准字
% * 不可定向标准字
% * 连通和与边界分支的恢复
% 
% 纲要内容结束

\section{标准模型的约化}
\label{g1:sec:0604}

上一节已把紧曲面的局部拼接信息压缩到有限条边上，但边数少并不意味着字已经容易辨认。例如，同一条边的两个字母相隔很远，可能隐藏着一个交叉帽；两个方向相反的边对互相穿插，则可能隐藏着一个手柄。本节用实际的切割和粘合把这些部件逐次取出。我们先处理闭曲面，最后把封住的边界圆盘恢复。这里要得到的是标准模型的存在；不同参数是否可能给出同胚曲面，由下一节另行判别。

\subsection{手柄和交叉帽的抽取}
\label{g1:subsec:0604:01}

先把多边形换成一种更容易局部操作的表示。\cref{g1:prop:one-vertex-band-model}把闭连通曲面写成一个顶点圆盘 \(D\)、有限条矩形带和一个面圆盘 \(F\) 的并：每条带的两条短边粘到 \(\partial D\) 上互不相交的小闭区间，且
\[
	B=D\cup\{\text{全部矩形带}\},\qquad \partial B\cong S^1,
\]
面圆盘 \(F\) 恰好封住这条边界。带的长边和顶点盘边界上没有被占用的弧组成 \(\partial B\)。这比抽象边字多保留了一点宽度，因而可以直接看见切下的对象是不是曲面、其边界有几条。

给 \(D\) 选一个定向。沿一条带把这个定向从一个粘合区间传到另一个区间，若与 \(D\) 原来的定向一致，称这条带无扭；若相反，称它有扭。此处说的扭转只记录两端的粘合方式，不需要把带嵌入三维空间再数半圈。顶点盘加一条有扭带是\term[mobiusdai]{Möbius 带}[Möbius band]\footnote{August Ferdinand Möbius（1790--1868），德国数学家、天文学家，研究射影几何、重心坐标及单侧曲面。}，可写成
\[
	\mathcal M=([0,1]\times[-1,1])/((0,t)\sim(1,-t)).
\]
顶点盘加一条无扭带则是环带。若两条带的四个粘合区间沿 \(\partial D\) 依次为 \(a,b,a,b\)，称这两个边对交错；这只说明端点在顶点盘周围的循环次序。

\begin{definition}[连通和与标准模型]\label{g1:def:surface-standard-models}
	从两个连通闭曲面 \(M,N\) 各去掉一个嵌入闭圆盘的内部，再沿两条剩下的边界圆周作同胚粘合，所得曲面称为它们的\term[liantonghe]{连通和}[connected sum]，记为 \(M\#N\)。本节先对所选圆盘及粘合实施这个构造，随后说明标准模型中这些选择为何不影响同胚类型。记
	\[
		\Sigma_{0,0}=S^2,\qquad \Sigma_{g,0}=\underbrace{T^2\#\cdots\#T^2}_{g\text{ 个}}\ (g\ge1),\qquad N_{k,0}=\underbrace{\mathbb RP^2\#\cdots\#\mathbb RP^2}_{k\text{ 个}}\ (k\ge1).
	\]
	再从相应闭模型去掉 \(b\) 个互不相交嵌入闭圆盘的内部，所得模型记为 \(\Sigma_{g,b}\) 和 \(N_{k,b}\)，其中 \(b\ge0\)。从一个闭曲面接上单孔环面，相应部件称为\term[shoubing]{手柄}[handle]；接上 Möbius 带，相应部件称为\term[jiaochamao]{交叉帽}[crosscap]。后者也就是与 \(\mathbb RP^2\) 作连通和，因为 \(\mathbb RP^2\) 去掉一个面内圆盘后正是 Möbius 带。
\end{definition}

连通和的两条边界各有环带邻域，把它们粘合后，接缝处的两个半圆盘组成一个圆盘，所以构造确实给出无边界曲面。\cref{g1:exm:surface-presentations}已分别从方形和圆盘对径识别得到环面与实射影平面；取其顶点及边的小带状邻域，就得到这里的单孔环面与 Möbius 带模型。下面把这个识别和余下部分同时说清楚。

\begin{lemma}[带条的抽取]\label{g1:lem:band-extraction}
	设闭连通曲面 \(M\) 由一个顶点闭盘 \(D\)、\(m\) 条矩形带及一个面闭盘 \(F\) 组成：各带的两个短边粘到 \(\partial D\) 上互不相交的区间，\(D\) 与全部带的并只有一个边界圆周，再由 \(F\) 封住。给 \(D\) 选定向，沿带传递此定向后，若与另一端 \(D\) 的定向一致，称该带无扭，否则称其有扭。两条分别标记为 \(a,b\) 的带称为交错，是指它们的粘合区间沿 \(\partial D\) 的循环次序为 \(a,b,a,b\)。若存在一条有扭带，则 \(M\cong\mathbb RP^2\#M'\)，其中 \(M'\) 有 \(m-1\) 条带的一顶点一面表示。若有两条交错的无扭带，则 \(M\cong T^2\#M'\)，其中 \(M'\) 有 \(m-2\) 条带的一顶点一面表示。
\end{lemma}
\begin{proof}
	第一种情形令 \(N\) 为 \(D\) 加选中的一条有扭带，第二种情形令 \(N\) 为 \(D\) 加选中的两条无扭带。未选带的粘合区间仍在 \(\partial N\) 上，暂时不将它们粘上。

	先核对 \(N\) 的类型。带条模型由粘合区间在顶点盘边界上的循环次序、每条带的两端方向决定：保持这些数据的边界同胚，可以在每个间隙弧上延拓，再延拓到顶点盘；带上按矩形坐标延拓。这些同胚在粘合区间上相容，故给出整个带面的同胚。上面的矩形 Möbius 模型恰好具有单条有扭带的粘合数据；环面方形模型中两条生成边的小邻域则具有两条交错无扭带的粘合数据。沿这个邻域的边界，从 \(a\) 的一侧出发，依次经过 \(a,b,a^{-1},b^{-1}\) 四段；其边界是一条圆周，补上一张盘就是方形的对边识别。因此第二种 \(N\) 是单孔环面。两种情形的 \(\partial N\) 都只有一条圆周。

	现在用一个新圆盘 \(D'\) 代替 \(N\)，将所有未选带原封不动地粘到 \(\partial D'\) 上对应的区间，再保留原来的面圆盘 \(F\)。换盘时边界上的粘合区间、间隙弧及其次序全部保持，因此新带面边界仍是原来供 \(F\) 封住的一条圆周。得到的 \(M'\) 是闭连通曲面，且正好少了一条或两条带。这里连通性也可以直接从模型读出：新顶点盘与每条余带相接，面盘沿整个边界相接。

	给 \(N\) 的边界补一个盘，所得闭曲面分别为 \(\mathbb RP^2\) 或 \(T^2\)。把这个补盘及 \(M'\) 中的 \(D'\) 各自去掉，再按先前的边界对应粘合，恰好恢复 \(M\)。这正是所声称的连通和分解，同时给出了余下带条的具体表示。
\end{proof}

上述证明没有假定“去掉一个手柄后仍应有某个亏格”。余下曲面是通过换盘实际造出来的，带条数确实下降；这将保证约化在有限步后停止。若读者习惯看边字，交叉帽的抽取还可写成下面的切粘等式，它稍后也用来把混合模型化为统一形式。

\begin{lemma}[一个同向边对的切粘]\label{g1:lem:crosscap-word-move}
	设一个边配对多边形的循环边字为 \(aUaV\)，字母 \(a\) 不在 \(U,V\) 中出现，其余每个边标号仍按原来的方式配对。则通过一次切割和一次重粘，可将它改为
	\begin{equation}\label{g1:eq:crosscap-word-move}
		aUaV\ \sim\ xxU^{-1}V.
	\end{equation}
	其中 \(x\) 为新边标号，\(U^{-1}\) 表示倒转 \(U\) 的次序并反转每个字母的指数；符号 \(\sim\) 表示配对所得曲面同胚，而非只表示基本群同构。
\end{lemma}
\begin{proof}
	在多边形内部取一条对角线 \(d\)，连接两条 \(a\) 边的起点。沿它切开，得到边字分别为
	\[
		aUd^{-1},\qquad daV
	\]
	的两个盘；若 \(U\) 为空，可把新角点稍作分离，或者直接保留原来相邻的 \(aa\)。把第一张盘反向读作 \(dU^{-1}a^{-1}\)，再沿两张盘中标为 \(a\) 的边粘合。粘合两盘的一整条边仍得到一个盘；剩余外边界按次序为 \(ddU^{-1}V\)。旧 \(a\) 变成内部接缝，新 \(d\) 的两份取代它成为配对边，重标为 \(x\) 即得所需字。旧边的配对始终未变，故\cref{g1:prop:polygon-moves}保证这是商曲面的同胚。图\ref{g1:fig:crosscap-word-move}标出实际切线与两张盘的边字。
\end{proof}

\begin{figure}[htbp]
	\centering
	% !TeX root = ../Book1.tex
\begin{tikzpicture}[scale=1,>=stealth,font=\small]
	\draw[thick,->] (0,2)--node[above]{$a$}(2,2);
	\draw[thick,->] (2,2)--node[right]{$U$}(2,0);
	\draw[thick,->] (2,0)--node[below]{$a$}(0,0);
	\draw[thick,->] (0,0)--node[left]{$V$}(0,2);
	\draw[blue!70!black,dashed,->,thick] (0,2)--node[above right]{$d$}(2,0);
	\fill (0,2) circle (1.7pt);
	\fill (2,0) circle (1.7pt);
	\node at (1,-0.65) {$aUaV$};
	\draw[->,thick] (2.8,1)--node[above,align=center]{沿 $d$ 切开\\沿 $a$ 重粘}(4.7,1);
	\draw[blue!70!black,thick,->] (5.4,2)--node[above]{$d$}(7.4,2);
	\draw[blue!70!black,thick,->] (7.4,2)--node[right]{$d$}(7.4,0);
	\draw[thick,->] (7.4,0)--node[below]{$U^{-1}$}(5.4,0);
	\draw[thick,->] (5.4,0)--node[left]{$V$}(5.4,2);
	\node at (6.4,-0.65) {$ddU^{-1}V$};
\end{tikzpicture}

	\caption[同向边对的切粘]{蓝色对角线连接两条 \(a\) 边的起点。先切出 \(aUd^{-1}\) 与 \(daV\)，将前一盘反向读边，再沿旧 \(a\) 边粘合，新的外边界为 \(ddU^{-1}V\)。图中 \(U,V\) 可以分别代表多条边。}
	\label{g1:fig:crosscap-word-move}
\end{figure}

\subsection{可定向标准字}
\label{g1:subsec:0604:02}

为什么可定向性恰好让抽取只产生手柄？所需信息已经在带条的两端，只要把局部定向沿带延拓一次，就能判别。

\begin{lemma}[带条与定向]\label{g1:lem:band-orientation}
	设闭连通曲面 \(M\) 由一个顶点闭盘 \(D\)、有限条矩形带及一个面闭盘组成：带的短边粘到 \(\partial D\) 上互不相交的区间，顶点盘与带的并只有一条边界，再由面盘封住。给 \(D\) 选定向；若沿一条带传递此定向后与另一端 \(D\) 的定向一致，就称该带无扭。则 \(M\) 可定向，当且仅当每条带都无扭。
\end{lemma}
\begin{proof}
	若曲面有定向，将它限制到顶点盘和一条带，沿带延拓必须返回同一个局部定向，因此该带无扭。反过来，先给顶点盘定向，无扭条件恰好保证每条矩形带可选一个定向，使其两端同时与顶点盘相容。各带内部互不相交，于是这些局部定向给整个带面 \(B\) 定向。\(B\) 的边界只有一条圆周，给面圆盘选择使接缝两侧相容的定向，即延拓到闭曲面。这里使用的是沿边界粘合时两侧诱导的边界方向相反；一条连通边界上相容性的符号不会跳变，所以一处相容便沿整条边界相容。
\end{proof}

\begin{proposition}[可定向标准字]\label{g1:prop:orientable-standard-word}
	对非负整数 \(g\)，记 \(\Sigma_{g,0}\) 为球面添 \(g\) 个手柄所得的曲面，\(\Sigma_{0,0}=S^2\)。每个非空紧连通、可定向、无边界拓扑曲面同胚于某个 \(\Sigma_{g,0}\)，其中 \(g\ge0\)。当 \(g\ge1\) 时，它有一顶点一面边字
	\[
		a_1b_1a_1^{-1}b_1^{-1}\cdots a_gb_ga_g^{-1}b_g^{-1}.
	\]
	当 \(g=0\) 时模型为 \(S^2\)，可直接看成两张圆盘沿整条边界粘合。
\end{proposition}
\begin{proof}
	用\cref{g1:prop:one-vertex-band-model}取得有限带条模型；若没有带，则顶点盘与面盘沿边界粘合，所得为球面。否则所有带均无扭。我们先证明必有两条带的端点交错。

	假设没有交错的边对。在圆周上的有限配对区间中，选一对，使其某一侧的开弧所含端点数最小。若这条弧非空，则弧内任一端点的配对端点也在弧内，否则它与所选边对交错。于是弧内又有一对，其一侧所含端点更少，矛盾。因此有一条带的两个粘合区间相邻，中间没有其他带端。该带无扭，所以它的一条长边与这段不含端点的顶点盘边界弧构成 \(\partial B\) 的一整条圆周。它不经过任何其他带；该带另一条长边还属于另一条边界圆周。这与 \(\partial B\) 只有一个分支矛盾。

	故存在交错的两条无扭带，\cref{g1:lem:band-extraction}给出 \(M\cong T^2\#M'\)，并把带数减少二。余曲面仍可定向：沿切下部件的边界将其余部分的定向延拓到新顶点盘即可。反复进行，每步减少两条带，有限步后没有带，余下为球面。于是得到有限个环面的连通和；与球面作连通和只是把一个盘取出后换回另一个盘，不改变同胚类型。

	最后说明边字而不只命名模型。单孔环面的边界在其顶点与两条核心边附近读出 \(aba^{-1}b^{-1}\)。把若干单孔模型通过一个连接各顶点的小圆盘串联，沿该盘外围依次读各部件，就得到上述交换子串。等价地，在两个面盘内部各去一小圆盘，用一条弧把新边界接到原多边形的顶点，沿弧切开后，连通和的面字是两者边字的串接；这条弧的两份作为内部接缝相互消去。所有端点汇到同一个顶点，最后补上一张面盘。这个构造也直接说明串接所得对象就是所取的环面连通和。
\end{proof}

注意，算法停止的理由是带条数这个非负整数严格下降，不是预先知道某个 Euler 示性数必须等于 \(2-2g\)。后一个等式将在下一节从模型计算出来，承担的是区分模型的职责。

\subsection{不可定向标准字}
\label{g1:subsec:0604:03}

若模型中有扭带，先取出一个交叉帽，余下部分可能仍不可定向，也可能变得可定向。因此直接抽取所得通常是手柄与交叉帽的混合。真正需要的新步骤，是在已有一个交叉帽时把一个手柄换成两个额外的交叉帽。

\begin{lemma}[一个手柄与一个交叉帽]\label{g1:lem:handle-crosscap-trade}
	存在同胚
	\[
		T^2\#\mathbb RP^2\cong\mathbb RP^2\#\mathbb RP^2\#\mathbb RP^2.
	\]
\end{lemma}
\begin{proof}
	左边的边字可循环读作 \(ccaba^{-1}b^{-1}\)。在下面的等式链中，每个带“切粘”的箭头均应用\cref{g1:eq:crosscap-word-move}或其逆操作，每个带“循环”的箭头只改变读边的起点；\(x,y,z\) 每次都是新的边标号：
	\begin{align*}
		ccaba^{-1}b^{-1}
		&\ \overset{\text{切粘}}{\sim}\ cb^{-1}a^{-1}ca^{-1}b^{-1}\\
		&\ \overset{\text{循环}}{\sim}\ a^{-1}ca^{-1}b^{-1}cb^{-1}\\
		&\ \overset{\text{切粘}}{\sim}\ xxc^{-1}b^{-1}cb^{-1}\\
		&\ \overset{\text{循环}}{\sim}\ b^{-1}cb^{-1}xxc^{-1}\\
		&\ \overset{\text{切粘}}{\sim}\ yyc^{-1}xxc^{-1}\\
		&\ \overset{\text{循环}}{\sim}\ c^{-1}xxc^{-1}yy\\
		&\ \overset{\text{切粘}}{\sim}\ zzx^{-1}x^{-1}yy.
	\end{align*}
	为核对第一步，可在基本公式中取 \(U=b^{-1}a^{-1}\)、\(V=a^{-1}b^{-1}\)；则 \(U^{-1}V=aba^{-1}b^{-1}\)，所以使用逆操作恰好得到第一行右端。其后三次分别抽取同向的 \(a^{-1}\)、\(b^{-1}\)、\(c^{-1}\) 边对。最后把 \(x^{-1}\) 重标为 \(u\)，边字成为 \(zzuuyy\)。这是三个实射影平面模型的字串接，因而给出所需同胚。每一步都作用于盘及边的实际识别，不是从群展示相同推测曲面相同。
\end{proof}

\begin{proposition}[不可定向标准字]\label{g1:prop:nonorientable-standard-word}
	对正整数 \(k\)，记 \(N_{k,0}\) 为 \(k\) 个 \(\mathbb RP^2\) 的连通和。每个紧连通、不可定向、无边界拓扑曲面同胚于某个 \(N_{k,0}\)，其中 \(k\ge1\)，并有一顶点一面边字
	\[
		c_1c_1c_2c_2\cdots c_kc_k.
	\]
\end{proposition}
\begin{proof}
	取得一顶点一面带条模型。由\cref{g1:lem:band-orientation}，至少有一条有扭带，先按\cref{g1:lem:band-extraction}取出一个交叉帽。以后只要余模型还有有扭带，就再取出一个；若余模型的带均无扭，则按可定向情形的证明取出手柄。每次带条数至少减一，所以终止后得到
	\[
		\underbrace{T^2\#\cdots\#T^2}_{h\text{ 个}}\#\underbrace{\mathbb RP^2\#\cdots\#\mathbb RP^2}_{r\text{ 个}},\qquad r\ge1.
	\]
	连通和中的部件可以重新排列、重新加括号：取一个有若干圆孔的球面，把每个去盘部件分别接到一条边界上，便同时表示这些次序与括号。交换两个孔可在平面圆盘里交换两小圆盘，再保持其外侧不动；边界参数的保向改变可在环带内延拓。标准的环面、实射影平面各有反射对称，它在指定的小面盘上反转边界方向，因此粘合方向的另一选择也给出同胚模型。于是可以每次把一个手柄与已有的一个交叉帽放在一起，应用\cref{g1:lem:handle-crosscap-trade}。每用一次，手柄数减一、交叉帽数增二；最终得到 \(N_{2h+r,0}\)。

	实射影平面的一个面盘边字为 \(cc\)，将这些部件按前一命题末尾的串接构造连接，所得标准字为 \(c_1^2\cdots c_k^2\)。每个部件中顶点均相同，串接处又识别这些顶点，故这是一个顶点、\(k\) 条边和一张面的表示。
\end{proof}

\begin{example}[Klein 瓶的两个边字]\label{g1:exm:klein-two-crosscaps}
	通常将方形按 \(aba^{-1}b\) 配对得到的曲面称为\term[kleinping]{Klein 瓶}[Klein bottle]\footnote{Felix Christian Klein（1849--1925），德国数学家，以变换群组织几何的 Erlangen 纲领及函数论研究著称。}。循环读为 \(ba^{-1}ba\)，在\cref{g1:eq:crosscap-word-move}中取 \(U=a^{-1}\)、\(V=a\)，立即得到 \(xxaa\)。因此 Klein 瓶同胚于 \(N_{2,0}\)。这个计算说明：是否已经出现相邻的平方，不是判别有无交叉帽的条件；真正的数据是同一边对的方向和全部粘合次序。
\end{example}

\subsection{连通和与边界分支的恢复}
\label{g1:subsec:0604:04}

对带边界曲面，先封盘再约化很方便，但最后还必须说明：得到闭模型之后，原先补入的盘可能位置很不规则，能否把它们移到标准的互不相交小盘？下面给出所需的有限配置结论。

\begin{lemma}[有限个圆盘的位置]\label{g1:lem:disk-configuration}
	设 \(M\) 是连通拓扑曲面，\(D_1,\ldots,D_b\) 与 \(E_1,\ldots,E_b\) 是其内部的两组两两不相交的嵌入闭圆盘，且各圆盘边界在曲面内有双侧环带邻域。则存在 \(M\) 的自同胚，将第一组圆盘作为无序集合送到第二组。若需要保持圆盘的编号，也可按任意指定的对应实现。
\end{lemma}
\begin{proof}
	先说明两个局部操作。一个有双侧环带的闭圆盘连同稍厚的外环带仍是闭圆盘。在这个较大盘中，将原盘的边界参数化为某个同心圆，再作分段线性的径向伸缩，可把原盘缩到内部任意小的圆盘；伸缩在较大盘的边界固定，因而可在外面接恒等映射。这里用\cref{g1:lem:disk-extension}把边界参数延拓到盘内，外侧环带则按乘积参数延拓。

	其次，在一个坐标圆盘内，可以把一个足够小的圆盘移到邻近位置，同时固定坐标圆盘外部。取半径 \(r<R\) 的同心圆盘，小盘包含待移动的盘，大盘仍在坐标域内。令 \(\rho(x)=\max\{0,\min\{1,(R-|x|)/(R-r)\}\}\)，则 \(\rho\) 在小盘上等于一，在大盘外等于零，并且 \(|\rho(x)-\rho(y)|\le L|x-y|\)，其中 \(L=1/(R-r)\)。若向量 \(v\) 足够短，使 \(L|v|<1\)，则映射 \(x\mapsto x+\rho(x)v\) 在小盘上作平移、在大盘外保持不动。由
	\[
		|x+\rho(x)v-y-\rho(y)v|\ge(1-L|v|)|x-y|
	\]
	知其单射。对每个 \(z\)，方程 \(x=z-\rho(x)v\) 的右端为压缩映射，逐次迭代收敛到唯一解：相邻迭代差按同一小于一的比例衰减，故为 Cauchy 列，极限满足方程。于是该映射满射，逆映射由上式连续，确为所需自同胚。把平移分成足够短的若干步，就可沿位于有限坐标盘中的道路移动小盘。

	还需保证已经放好的盘不会挡住道路。连通曲面去掉有限个互不相交内盘的内部后仍连通：所得空间是有边界的局部道路连通空间，若有多个连通分支，每条新边界圆周完全落在其中一个分支；把所去的盘沿该圆周补回，只能与这个分支相接，不会连接不同分支。局部半圆盘及所补盘的邻域使这些分支补盘后仍给出原曲面的不交开分解，与连通性矛盾。因此余下曲面道路连通，利用边界环带可将道路略微推入其内部。

	现逐个安置圆盘。目标盘之外有非空开集：有限个互不相交的内闭盘不可能覆盖整个连通曲面。先将全部原盘缩得足够小，使该开集内有一个与它们不交的坐标小盘，再在这里选 \(b\) 个暂存小盘。将原盘逐个送到这些暂存位置，移动时保持其他当前圆盘不动；此阶段目标盘只标记最终位置，允许移动操作经过它们。随后逐个沿避开其他当前圆盘的道路移到指定目标盘内部，再在目标盘及其外环带中作径向扩张，填满目标盘。已经放好的目标盘保持不动。每一步只有有限个盘需要避开，前段保证可选择道路；每条道路的紧性又保证只需有限个局部移动。复合所得自同胚即实现所给对应。
\end{proof}

暂存位置的作用是避开一种实际困难：原盘可能横跨几个目标盘，直接逐个扩张到目标位置会撞上尚未搬走的盘。先全部移到目标区域之外，再逐个放入，便把这个障碍消除了。对于 \(b=0\)，结论就是恒等映射。

\begin{theorem}[紧曲面标准模型的存在]\label{g1:thm:surface-standard-existence}
	记 \(\Sigma_{g,b}\) 为球面添 \(g\) 个手柄后挖去 \(b\) 个互不相交闭圆盘的内部所得的曲面，\(N_{k,b}\) 为 \(k\) 个 \(\mathbb RP^2\) 的连通和后作同样的 \(b\) 次挖孔，其中 \(g,b\ge0\)、\(k\ge1\) 均为整数。每个非空紧连通、Hausdorff、第二可数的拓扑曲面 \(M\)，边界允许非空，都同胚于某个
	\[
		\Sigma_{g,b}\quad(g,b\ge0),\qquad\text{或}\qquad N_{k,b}\quad(k\ge1,\ b\ge0).
	\]
	其中 \(b\) 为 \(M\) 的边界分支数；可定向的 \(M\) 得到第一类，不可定向的 \(M\) 得到第二类。
\end{theorem}
\begin{proof}
	由\cref{g1:prop:cap-boundary}，\(\partial M\) 是有限条圆周的并，逐一封盘得到闭连通曲面 \(\widehat M\)。若 \(M\) 可定向，其定向在每个补盘上按边界相容的方式延拓；若 \(\widehat M\) 可定向，其限制当然给 \(M\) 定向。因此封盘不改变可定向性。

	由前面的两种标准字结论，\(\widehat M\) 同胚于 \(\Sigma_{g,0}\) 或 \(N_{k,0}\)。补入的 \(b\) 张盘在这个同胚下仍是两两不交的闭盘，接缝两侧的环带使它们满足\cref{g1:lem:disk-configuration}。用该引理把这些盘移到标准模型中指定的 \(b\) 个小盘，再去掉它们的内部，便得到 \(M\cong\Sigma_{g,b}\) 或 \(M\cong N_{k,b}\)。去盘后每张盘贡献一条边界且没有别的边界，所以参数 \(b\) 正好为原边界分支数。
\end{proof}

圆盘配置引理也补足了连通和构造的选择问题：改变取盘位置可由曲面自同胚实现；两个同向的圆周粘合参数之差，可以沿环带作同伦延拓。具体地，将保向圆周同胚提升为严格递增的 \(f:\mathbb R\to\mathbb R\)，满足 \(f(t+1)=f(t)+1\)，则 \((1-s)f(t)+st\) 仍严格递增，给出到恒等映射的同伦；把参数 \(s\) 取为环带的横坐标即可延拓。标准模型还有反射对称，允许反转某一粘合边界的方向。各部件接到带孔球面上的描述说明重新加括号和交换部件也不改变结果。因此 \(\Sigma_{g,b}\)、\(N_{k,b}\) 在固定参数下确有明确的同胚类型。

例如，\(\Sigma_{0,1}\) 是圆盘，\(\Sigma_{0,2}\) 是环带，\(N_{1,1}\) 是 Möbius 带；\(\Sigma_{0,3}\) 是三孔球面。若从 \(T^2\#T^2\#\mathbb RP^2\) 去掉三个内盘，连续使用两次手柄与交叉帽的换位得到 \(N_{5,3}\)。这些参数是由实际切粘读出的。下一节将证明它们不能再被改成另一组参数，从而把本节的模型存在性提升为分类定理。
% 纲要条目排为编号小节标题，并进入目录与 PDF 书签；说明段仍保留于纲要注释中。
